% !TeX root = ../main.tex
\section{数值实验}
\label{exp:section}

数值实验从换电路径与充电功率的联合优化、RL 终端价值和运行条件三个方面评价所提方法。本次完成第~\ref{exp:terminal}~节的零终端价值实验；RL 训练、RL 对照及完美信息结果尚未开展。各预测域使用相同的道路网络、七个测试日、预约与随机请求、误差、日前路径和初始库存，需求预测仅使用决策时刻可获得的信息。

道路包含 11 个换电站、6 个出入口和 26 个 O--D 对。节点名称与坐标来自 \path{poi_node_6e.csv}，站点与出入口里程分别来自 \path{distance_stations.csv}、\path{distance_6e_entrance.csv}；站点电池配置、电价和服务费分别来自 \path{11s_configuration_info.csv}、\path{11s_elec_price_info.csv}、\path{11s_service_fee_info.csv}。O--D 与人口数据来自 \path{subpaths_6e.csv} 和 \path{population_6e.csv}。预约 O--D 抽样权重为起终点人口乘积除以距离平方，入场小时权重来自 \path{flowVolume.csv}，小时内入场时刻均匀抽样。

根据 \texttt{demand\_info.csv} 中站点、小时和星期对应的真实需求均值，独立进行泊松抽样，生成连续十周、共 70 天的\textbf{基于真实均值的泊松合成场景}。从每种星期对应的十天中随机选一天，周一至周日的测试日编号依次为 1、37、3、4、19、13、14；其余 63 天作为后续训练和验证候选池，尚未进一步划分。每个测试日均生成 200 名预约用户和 200 个随机请求。随机请求以该日泊松样本的站点与小时联合频数为权重进行固定总量的多项分配，再在小时内均匀抽样到达时刻；零总量样本报错，不更换测试日。随机需求预测使用原文件中对应星期的均值，归一化至每日期望 200，并按累计期望量生成确定性预测请求，不读取测试日未来实际请求。数据来源、随机流及场景哈希保存于 \texttt{data/final\_real\_data\_v1/manifest.json}；各随机流均由基础种子 1 按日期和用途派生。

计算环境见表~\ref{exp:environment}，主要参数见表~\ref{exp:parameters}。本章实验取 $\delta=15$~min、$K=2$，覆盖模型说明中默认的 5~min 和 $K=3$：充电功率每段更新，换电路径每 30~min 允许更新。$H=4,5,6,7,8$ 分别对应 60、75、90、105、120~min，$H$ 本身即时间段数。单次求解时限为 120~s，目标相对 gap 为 $10^{-4}$（0.01\%），每个日任务最多使用 16 线程，各任务串行运行。达到时限但存在可行解时执行该解，并保留实际 gap；没有可行解时停止任务并保存诊断。

\begin{table}[!htbp]
  \centering
  \caption{计算环境}
  \label{exp:environment}
  \small
  \begin{tabularx}{0.88\linewidth}{@{}lX@{}}
    \toprule
    项目 & 配置 \\
    \midrule
    CPU & Intel Core i7-12650H，10 核、16 逻辑线程 \\
    系统可见内存 & 15.7~GiB \\
    操作系统 & Windows，系统构建号 10.0.26200 \\
    Python 版本 & 3.10.19 \\
    COPT 版本 & 8.0.6 \\
    GPU / PyTorch & 本次零终端价值实验未使用 \\
    \bottomrule
  \end{tabularx}
\end{table}

\begin{table}[!htbp]
  \centering
  \caption{主要实验参数}
  \label{exp:parameters}
  \small
  \setlength{\tabcolsep}{5pt}
  \begin{tabularx}{\linewidth}{@{}Xll@{}}
    \toprule
    参数 & 符号 & 设置 \\
    \midrule
    站点数量 & $|\mathcal S|$ & 11 \\
    出入口与 O--D 数量 & --- & 6，26 \\
    电池总数与初始状态 & --- & 250 块，全部满电 \\
    电池容量 & $E_B$ & 100~kWh \\
    各站充电效率 & $\eta_i$ & 95\% \\
    设备功率上限 & $\overline P_{i,b}^{\mathrm{ch}}$ & 60~kW \\
    站点总充电功率上限 & $\overline P_i^{\mathrm{sta}}$ & 960~kW \\
    用户最长等待时间 & --- & 30~min \\
    路径调整与预约失败成本 & $c^{\mathrm{adj}},c^{\mathrm{fail}}$ & 10 元/次，1000 元/人 \\
    充电更新周期 & $\delta$ & 15~min \\
    路径更新周期 & $K\delta$ & 30~min \\
    预测域长度 & $H$ & 4、5、6、7、8 段 \\
    需求期长度 & $N$ & 96 段，24~h \\
    日需求量：预约 / 随机 & --- & 200 / 200 \\
    车速、满电续航 & --- & 75~km/h，300~km \\
    出口 SOC 下限、间距剪枝阈值 & --- & 0.1，100~km \\
    本次终端价值权重 & $\beta$ & 0 \\
    候选池 / 测试日数量 & --- & 63 / 7 \\
    单次求解时限、目标相对 gap & --- & 120~s，0.01\% \\
    基础随机种子 & --- & 1 \\
    \bottomrule
  \end{tabularx}
\end{table}

按 $s_1$ 至 $s_{11}$ 顺序，各站电池数为 $[22,22,22,18,20,25,23,20,28,28,22]$；变量、执行和统计均按实际槽数处理。预约实际入口 SOC 在 $(0.5,1]$ 上均匀生成，上报 SOC 限于 $[0.5,1]$；上报时间误差上限为 10~min，SOC 误差上限为 0.05，在各自可行区间内均匀生成。物理路段行驶时间乘数独立取自 $U[0.9,1.1]$；随机退回 SOC 取自 $U[0.1,0.3]$，预测值取 0.2。为保证入口误差下的首段可达性，日前及未入场用户的候选网络使用 $\max(0.5,s^{\mathrm{report}}-0.05)$ 作为可达性计算的入口 SOC；预测退回电量使用报告值，实际执行使用真实值。100~km 沿用第~\ref{sec:network}~节的剪枝规则，保留必要可行路径和既有计划，并非对所有弧施加硬下限。

24:00 后停止新增预约与随机需求，随机预测同样不再生成新请求；预测窗口保持指定 $H$，已有用户继续处理至全部完成服务或失败。预约完成按其剩余换电安排为空且可满足出口 SOC 要求判定，最后换电后的无服务行驶不再延长仿真。价格按 24 小时循环，跨日收入和费用均归入原需求日，不计库存残值、不强制补满。按最长行程、行驶误差和逐站最大等待推导的安全上限为 186 段（46.5~h）；到达该上限仍有待处理用户时标记异常。需求期、实际执行期和跨日处理期分别记录。

运营净收益为实际服务收入减去充电、路径调整和预约失败成本。换电收入按实际补充电量乘以当时电价与服务费之和计算，充电成本按电网输入电量另行扣除。预约失败率为失败用户数除以 200，同一用户仅罚一次；随机服务率为服务请求数除以 200。平均等待时间先在每个日场景内对全部已服务换电请求计算，再对七天取平均。路径调整次数仅记录向在途用户实际发布的不同换电站序列，未入场计划变化不计实际费用。

所有主指标均报告固定七个测试日的算术均值；每个 $H$ 每天仅运行一次，共 35 个日任务，属于同一基础种子下的跨日比较，不是多种子重复实验。表~\ref{exp:value_results} 的求解均值及 P95 先在每日需求期的 96 轮内计算，再对七天取平均；含跨日阶段的全部轮次统计另行报告。P95 为日内第 95 百分位数，七天不按轮次数加权。样本分母为零时记为“不适用”。



\FloatBarrier
\subsection{The role of RL terminal value function}
\label{exp:terminal}

本节拟比较 $\beta=0$ 的无终端价值设置与 $\beta=1$ 的 RL 终端价值设置，预测域取 $H=4,\ldots,8$。当前已完成零终端价值的 35 个日场景；其结果刻画短预测域下的运营表现和计算代价，尚不能单独证明 RL 终端价值的收益。

终端价值网络的训练设置保留在表~\ref{exp:training}，待实际训练后填写。表~\ref{exp:value_results} 按 $H$ 列出已完成的零终端价值结果，RL 及完美信息对照保持待填。

\begin{table}[!htbp]
  \centering
  \caption{终端价值网络的训练设置}
  \label{exp:training}
  \small
  \begin{tabularx}{\linewidth}{@{}Xc@{}}
    \toprule
    项目 & 设置 \\
    \midrule
    输入维度 & 待填 \\
    隐藏单元数 & 待填 \\
    参数总数 & 待填 \\
    训练场景数量 & 待填 \\
    采样批次数与每批轨迹数 & 待填 \\
    每批参数更新次数 & 待填 \\
    回归批量大小 & 待填 \\
    优化器与学习率 & 待填 \\
    独立训练次数 & 待填 \\
    验证频率与网络选择规则 & 待填 \\
    实际训练耗时 & 待填 \\
    \bottomrule
  \end{tabularx}
\end{table}

\begin{table}[!htbp]
  \centering
  \caption{终端价值实验的运营结果（固定七日算术均值）}
  \label{exp:value_results}
  \small
  \setlength{\tabcolsep}{3pt}
  \begin{tabularx}{\linewidth}{@{}Xcccccc@{}}
    \toprule
    方法 & \shortstack{净收益\\元} & \shortstack{预约\\失败率/\%} & \shortstack{随机用户\\服务率/\%} & \shortstack{等待时间\\min} & \shortstack{调整次数\\次} & \shortstack{求解时间\\均值/P95，s} \\
    \midrule
    零终端 $H=4$ & $-96706.58$ & 60.57 & 62.36 & 11.94 & 45.71 & 0.11/0.23 \\
    零终端 $H=5$ & $-83734.73$ & 55.00 & 64.07 & 11.53 & 54.00 & 0.18/0.42 \\
    零终端 $H=6$ & $-49887.07$ & 40.07 & 70.29 & 12.10 & 76.00 & 1.13/2.06 \\
    零终端 $H=7$ & 37180.19 & 0.21 & 86.93 & 12.87 & 141.14 & 23.09/120.01 \\
    零终端 $H=8$ & 37639.56 & 0.43 & 89.57 & 12.55 & 198.71 & 50.15/120.19 \\
    \midrule
    RL $H=4$ & 待填 & 待填 & 待填 & 待填 & 待填 & 待填 \\
    RL $H=5$ & 待填 & 待填 & 待填 & 待填 & 待填 & 待填 \\
    RL $H=6$ & 待填 & 待填 & 待填 & 待填 & 待填 & 待填 \\
    RL $H=7$ & 待填 & 待填 & 待填 & 待填 & 待填 & 待填 \\
    RL $H=8$ & 待填 & 待填 & 待填 & 待填 & 待填 & 待填 \\
    完美信息上界 & 待填 & 待填 & 待填 & 待填 & 待填 & 待填 \\
    \bottomrule
  \end{tabularx}
\end{table}

\begin{table}[!htbp]
  \centering
  \caption{零终端价值的求解与跨日处理诊断}
  \label{exp:zero_diagnostics}
  \small
  \setlength{\tabcolsep}{3pt}
  \begin{tabularx}{\linewidth}{@{}cXXXXXX@{}}
    \toprule
    $H$ & \shortstack{限时轮次\\比例/\%} & \shortstack{gap 达标\\比例/\%} & \shortstack{平均实际\\gap/\%} & \shortstack{跨日时长\\h} & \shortstack{终态库存\\kWh} & \shortstack{终态满电\\电池数} \\
    \midrule
    4 & 0.00 & 100.00 & 0.000052 & 3.11 & 7054.81 & 12.29 \\
    5 & 0.00 & 100.00 & 0.000163 & 5.00 & 6609.73 & 12.00 \\
    6 & 0.24 & 99.76 & 0.000899 & 5.82 & 5512.72 & 6.00 \\
    7 & 10.19 & 89.81 & 0.163012 & 11.18 & 4864.18 & 0.71 \\
    8 & 26.62 & 73.38 & 0.969465 & 11.57 & 4735.21 & 1.14 \\
    \bottomrule
  \end{tabularx}
  \par\smallskip
  \parbox{\linewidth}{\footnotesize 注：各项先按日统计，再对七日等权平均。此表的求解比例及 gap 使用包含跨日处理的全部轮次；gap 达标指实际相对 gap 不超过 0.01\%。终态满电数量为日终计数的均值，故可为小数。}
\end{table}

当 $H$ 从 4 增至 7 时，七日平均净收益由 $-96706.58$ 元增至 37180.19 元，预约失败率由 60.57\% 降至 0.21\%，随机服务率由 62.36\% 增至 86.93\%。$H=4$ 时每日平均充电费用仅为 662.48 元，而预约失败费用为 121142.86 元。以随机请求的预测退回 SOC 0.2 为例，在 60~kW、95\% 效率下将 100~kWh 电池充满需要约 84.21~min，即至少六个完整充电时段，之后的边界才能再次服务。较短窗口难以同时看到这笔充电成本及其后续服务收益；$H=7$ 可覆盖六段充电后的服务边界。这为收益在 $H=6$ 与 $H=7$ 之间明显变化提供了物理解释，但不构成对所有请求或场景的阈值保证。

$H=8$ 相比 $H=7$ 的七日平均净收益增加 459.37 元（约 1.24\%），随机服务率提高 2.64 个百分点，需求期单轮平均求解时间则由 23.09~s 增至 50.15~s。预约失败总人数由七天合计 3 人增至 6 人，因此并非所有服务指标均随 $H$ 单调改善。需求期每轮二元变量数的七日均值由 $H=4$ 的约 2489 个增至 $H=8$ 的约 3933 个；限时求解比例也明显上升。含跨日阶段时，$H=4,\ldots,8$ 的单轮求解均值/P95 七日均值依次为 0.10/0.22、0.15/0.40、0.90/1.58、15.79/120.00、34.08/120.15~s，低于仅需求期统计的相应均值。

表~\ref{exp:zero_diagnostics} 表明，0.01\% 是求解目标，不能表述为全部轮次均达到的精度。$H=7$ 和 $H=8$ 全部日场景中的最大实际 gap 分别为 6.02\% 和 16.97\%。这些有限时长可行解对应的运营轨迹均完成核验，结果比较仍受求解预算影响，不能将差异全部归因于预测域长度。求解器退出和计时开销使少数记录略高于 120~s。

35 条轨迹共包含 4387 个执行时段。独立重建账务及指标后，每日预约完成数加失败数均为 200，随机服务数加超时数均为 200，最终无未处理用户。满电交付、电池 SOC 与功率、预约优先及同类先到先服务、后站到达对前站服务的依赖、等待截止边界及实际发布计费均通过核验。跨日处理最长为 14.5~h，低于推导的安全上限；道路最长 O--D 为 1140~km，长行程与晚间入场可产生较长跨日服务。当前结论仅适用于初始全满且无库存残值结算的独立需求日，不能直接解释为连续多日稳态收益。

\begin{figure}[!htbp]
  \centering
  \fbox{\begin{minipage}[c][32mm][c]{0.9\linewidth}
    \centering
    待填：价值网络训练过程及终端价值消融结果\\[4pt]
    横轴与单位、RL 对照曲线及按测试日统计的差异待填。
  \end{minipage}}
  \caption{终端价值学习过程及运营表现（待填）}
  \label{exp:value_figure}
\end{figure}

\todoresult{待填：完成 RL 训练及配对测试后，比较相同预测域下的各类终端价值；完美信息结果另行计算，不从零终端价值结果推断。}


\FloatBarrier
\subsection{The value of route and charging power optimization}
\label{exp:joint}

为考察两项决策的作用，设置表~\ref{exp:joint_methods}所列四组方法。各组终端价值均取零，使用相同预测域长度、需求输入和服务规则。服务时间与满电电池匹配均由优化模型确定。采用日前路径的组别在运营过程中沿用日前计算的换电站序列；优化换电路径的组别按规定的更新周期调整在途预约用户的换电安排。

\begin{table}[!htbp]
  \centering
  \caption{换电路径与充电功率联合优化的对照设置}
  \label{exp:joint_methods}
  \small
  \begin{tabularx}{\linewidth}{@{}Xlll@{}}
    \toprule
    方法 & 换电路径 & 充电功率 & 终端价值 \\
    \midrule
    基准方法 & 日前路径 & 基准充电规则 & 0 \\
    仅优化换电路径 & 在线优化 & 基准充电规则 & 0 \\
    仅优化充电功率 & 日前路径 & 在线优化 & 0 \\
    联合优化 & 在线优化 & 在线优化 & 0 \\
    完美信息上界 & 在线优化 & 在线优化 & 0 \\
    \bottomrule
  \end{tabularx}
\end{table}

基准充电规则将站点总充电功率上限按设备数量等分，实际功率同时满足设备上限和电池可充电量限制：
\begin{equation}
 P_{i,b,n}^{\mathrm{base}}=
 \min\left\{
 \overline P_{i,b}^{\mathrm{ch}},\;
 \frac{\overline P_i^{\mathrm{sta}}}{|\mathcal B_i|},\;
 \frac{E_B(1-S_{i,b,n}^{+})}{\eta_i\delta}
 \right\}.
 \label{eq:baseline_power}
\end{equation}
该规则使用本时刻完成换电后的 SOC。每台设备保留等分得到的功率额度，满电电池对应设备的功率为零。采用基准规则的组别在预测域内施加 $P_{i,b,n}=(1-d_n^{\mathrm{end}})P_{i,b,n}^{\mathrm{base}}$，使各运营时间段使用基准功率，结束后的时间段功率为零。其分段线性表达见附录。式中 $\delta$ 以小时计，因而第三项的单位为 kW。

表~\ref{exp:joint_results}记录运营指标，表~\ref{exp:joint_costs}记录收入与成本分解。通过基准方法与两组单项优化的比较，考察路径调整和充电安排各自的作用；通过两组单项优化与联合优化的比较，考察同时优化两项决策的运营表现。各项比较使用相同测试场景的配对结果。

\begin{table}[!htbp]
  \centering
  \caption{联合优化消融实验的运营结果}
  \label{exp:joint_results}
  \small
  \setlength{\tabcolsep}{3pt}
  \begin{tabularx}{\linewidth}{@{}Xcccccc@{}}
    \toprule
    方法 & \shortstack{净收益\\元} & \shortstack{预约\\失败率/\%} & \shortstack{随机用户\\服务率/\%} & \shortstack{等待时间\\min} & \shortstack{调整次数\\次} & \shortstack{求解时间\\均值/P95，s} \\
    \midrule
    基准方法 & 待填 & 待填 & 待填 & 待填 & 待填 & 待填 \\
    仅优化换电路径 & 待填 & 待填 & 待填 & 待填 & 待填 & 待填 \\
    仅优化充电功率 & 待填 & 待填 & 待填 & 待填 & 待填 & 待填 \\
    联合优化 & 待填 & 待填 & 待填 & 待填 & 待填 & 待填 \\
    完美信息上界 & 待填 & 待填 & 待填 & 待填 & 待填 & 待填 \\
    \bottomrule
  \end{tabularx}
\end{table}

\begin{table}[!htbp]
  \centering
  \caption{联合优化消融实验的收入与成本分解（元）}
  \label{exp:joint_costs}
  \small
  \begin{tabularx}{\linewidth}{@{}Xcccc@{}}
    \toprule
    方法 & 服务收入 & 充电成本 & 路径调整成本 & 预约失败成本 \\
    \midrule
    基准方法 & 待填 & 待填 & 待填 & 待填 \\
    仅优化换电路径 & 待填 & 待填 & 待填 & 待填 \\
    仅优化充电功率 & 待填 & 待填 & 待填 & 待填 \\
    联合优化 & 待填 & 待填 & 待填 & 待填 \\
    完美信息上界 & 待填 & 待填 & 待填 & 待填 \\
    \bottomrule
  \end{tabularx}
\end{table}

\todoresult{待填：根据运营指标和成本分解，分析换电路径优化、充电功率优化及二者联合优化的作用，并给出重复实验的统计结果。}


\FloatBarrier
\subsection{敏感性分析}
\label{exp:sensitivity}

敏感性分析考察需求强度、预约和随机人数的比例、电池库存配置、用户旅行时间的波动性。表~\ref{exp:sensitivity_settings}记录各因素的基准值和变化范围。

\begin{table}[!htbp]
  \centering
  \caption{敏感性分析的参数设置}
  \label{exp:sensitivity_settings}
  \small
  \begin{tabularx}{\linewidth}{@{}Xlll@{}}
    \toprule
    因素 & 基准值 & 变化范围 & 取值数量 \\
    \midrule
    需求强度 & 待填 & 待填 & 待填 \\
    预约与随机用户比例 & 待填 & 待填 & 待填 \\
    各站电池数量 $|\mathcal B_i|$ & 待填 & 待填 & 待填 \\
    用户旅行时间波动性 & 待填 & 待填 & 待填 \\
    \bottomrule
  \end{tabularx}
\end{table}



\begin{table}[!htbp]
  \centering
  \caption{敏感性分析结果（各参数取值分别填列）}
  \label{exp:sensitivity_results}
  \small
  \setlength{\tabcolsep}{3pt}
  \begin{tabularx}{\linewidth}{@{}Xcccccc@{}}
    \toprule
    因素、取值及方法 & \shortstack{净收益\\元} & \shortstack{预约\\失败率/\%} & \shortstack{随机用户\\服务率/\%} & \shortstack{等待时间\\min} & \shortstack{调整次数\\次} & \shortstack{求解时间\\均值/P95，s} \\
    \midrule
    需求强度：待填 & 待填 & 待填 & 待填 & 待填 & 待填 & 待填 \\
    预约与随机用户比例：待填 & 待填 & 待填 & 待填 & 待填 & 待填 & 待填 \\
    电池库存配置：待填 & 待填 & 待填 & 待填 & 待填 & 待填 & 待填 \\
    用户旅行时间波动性：待填 & 待填 & 待填 & 待填 & 待填 & 待填 & 待填 \\
    \bottomrule
  \end{tabularx}
\end{table}

\begin{figure}[!htbp]
  \centering
  \fbox{\begin{minipage}[c][36mm][c]{0.9\linewidth}
    \centering
    待填：预测域长度、需求强度、电池库存和站点总充电功率上限的敏感性结果\\[4pt]
    各参数取值、运营指标曲线及求解时间曲线待填。
  \end{minipage}}
  \caption{不同运行条件下的运营表现与求解时间（待填）}
  \label{exp:sensitivity_figure}
\end{figure}

\todoresult{待填：分别分析预测域长度、需求强度、电池数量和站点总充电功率上限变化时的收益、服务表现及求解时间，并说明各项结论所对应的参数范围。}

